\textbf{结论名称：} 对角互补模型

\textbf{适用条件：} 四边形中有一组对角互补；若需旋转全等，还需一组邻边相等。

\textbf{变量说明：} 四边形 \(ABCD\)，内角 \(\angle A,\angle B,\angle C,\angle D\)；旋转中心常取相等邻边的公共顶点（如点 \(A\)）。

\textbf{核心公式：}
(1) 对角互补 \(\Rightarrow\) 四点共圆
\[
\boxed{\angle A + \angle C = 180^\circ \;\Longrightarrow\; A,B,C,D \text{ 四点共圆}}
\]
(2) 对角互补且邻边相等 \(\Rightarrow\) 旋转全等
\[
\boxed{\angle B + \angle D = 180^\circ,\; AB = AD \;\Longrightarrow\; \triangle ABC \cong \triangle ADE}
\]
（将 \(\triangle ABC\) 绕点 \(A\) 旋转，使 \(AB\) 与 \(AD\) 重合，\(E\) 落在射线 \(CD\) 上）

\textbf{使用场景：} 四点共圆判定；与角平分线、截长补短结合，将分散的边角集中以解决线段和差问题。

\textbf{简单例子：} 在四边形 \(ABCD\) 中，\(\angle A = 105^\circ\)，\(\angle C = 75^\circ\)。
求证 \(A,B,C,D\) 四点共圆，并求 \(\angle B + \angle D\) 的度数。
解：\(\angle A + \angle C = 105^\circ + 75^\circ = 180^\circ\)，由对角互补得四点共圆。在圆内接四边形中，对角互补，故 \(\angle B + \angle D = 180^\circ\)。